\section{Zielsetzung}\label{zielsetzung}

Das übergeordnete Ziel dieser Arbeit ist die Entwicklung und Evaluation eines erklärbaren ML-basierten Intrusion Detection Systems für IoT-Netzwerke, das auf dem CICIoT2023-Datensatz trainiert und bewertet wird. Das System soll zwei Eigenschaften vereinen, die in der bestehenden Literatur bislang nicht gemeinsam auf diesem Datensatz untersucht wurden: eine hohe Erkennungsleistung gemessen an der Balanced Accuracy sowie interpretierbare Entscheidungsgrundlagen durch SHAP.

Daraus ergeben sich drei konkrete wissenschaftliche Beiträge:

Der erste Beitrag besteht in der erstmaligen Integration von SHAP in ein baumbasiertes Klassifikationsexperiment auf dem CICIoT2023-Datensatz, evaluiert anhand der Balanced Accuracy als primärer Bewertungsmetrik. Diese Kombination existiert in der Literatur bisher nicht und stellt eine direkt belegbare Forschungslücke dar.

Der zweite Beitrag besteht in einer quantitativen Ablation Study, die den Einfluss der PCA-Dimensionsreduktion auf Modellleistung und SHAP-Feature-Importance systematisch untersucht. Diese Analyse fehlt im Ausgangspaper vollständig.

Der dritte Beitrag besteht in einer semantischen Validierung der SHAP-Erklärungen. Die identifizierten Feature-Importance-Ranglisten werden mit dem fachlichen Wissen über die Netzwerksignaturen der jeweiligen Angriffstypen verglichen, um zu prüfen, ob das Modell tatsächlich sicherheitsrelevante Merkmale lernt oder statistische Artefakte des Datensatzes.

Die drei Forschungsfragen, die diese Ziele operationalisieren, lauten wie folgt:

Kann ein Gradient-Boosting-Klassifikator auf dem CICIoT2023-Datensatz eine hohe Balanced Accuracy bei der Erkennung von Cyberangriffen erzielen, und lassen sich seine Entscheidungen durch SHAP so erklären, dass die relevanten Netzwerkmerkmale pro Angriffsklasse für Sicherheitsanalysten nachvollziehbar werden?

Inwiefern verändert sich die Modellbewertung, wenn statt der einfachen Accuracy die Balanced Accuracy als primäre Metrik eingesetzt wird, und welche Angriffsklassen werden durch das Klassenungleichgewicht systematisch schlechter erkannt?

Welchen messbaren Einfluss hat die PCA-basierte Dimensionsreduktion auf die Balanced Accuracy und die SHAP-Erklärungen, und welche Netzwerkmerkmale sind gemäß SHAP für die Klassifikation der einzelnen Angriffstypen entscheidend?
